%!TEX root = ../main.tex
\section{Feedback-Driven Data Evolution}
\label{sec:data-engines}

A static pipeline applies a fixed sequence of operations to an available pool. A data engine differs because evidence from the current learner changes which records are collected, generated, labeled, weighted, or checked in a later round. Repeating this loop produces the data self-evolution defined in Section~\ref{sec:introduction}. The object that evolves is not only a dataset snapshot. It is the training distribution together with its units, roles, weights, and preparation rules.

This definition excludes two nearby practices. Re-running the same cleaning script is workflow automation. Periodically fine-tuning on every new log is a model-update schedule. Neither closes a data loop unless the behavior of the current learner determines what is sought, changed, retained, or assigned to a new objective.

\subsection{What Closes the Loop}

A complete round contains four functions. The current model is evaluated on a declared target. Its outcomes, interventions, uncertainty, and coverage gaps are converted into structured evidence. A controller then chooses a data operation within human, compute, simulation, and robot budgets. A successor is trained and deployed only after a controlled comparison with the current model. Data, verifier, and model versions must share lineage because changing any one of them changes the evidence observed in the next round.

Table~\ref{tab:data-engines} compares representative loops through their feedback, next data decision, and safeguard. RoboCat uses practice to extend the training mixture, while AutoRT uses foundation models to propose feasible fleet tasks and collect experience~\cite{bousmalis2023robocat_2306_11706,ahn2024autort_2401_12963}. DexFlyWheel performs a simulation-centered loop seeded by a human demonstration~\cite{zhu2025dexflywheel_2509_23829}. The $\pi^{*}_{0.6}$ report names its experience-learning loop RECAP~\cite{intelligence20250_2511_14759}. Robot-Powered Data Flywheels adapts an upstream perception model from observations encountered during deployment, showing that the updated learner need not be the action policy~\cite{grannen2025robot_2511_19647}. These systems share a feedback pattern, but the semantics of their new data and the evidence used to deploy a successor are different.

\begin{table*}[t]
\centering
\caption{Data engines compared as feedback controllers.  The discriminating
questions are what detects a gap, which data action follows, how the learner is
updated, and what prevents a locally improving loop from drifting or collapsing.}
\label{tab:data-engines}
\scriptsize
\setlength{\tabcolsep}{3.5pt}
\begin{tabularx}{\textwidth}{@{}p{2.8cm}p{3.5cm}p{6.15cm}X@{}}
\toprule
System & Gap/feedback signal & Data action and model update & Control mechanism and remaining limit \\
\midrule
RoboCat / AutoRT
~\cite{bousmalis2023robocat_2306_11706,ahn2024autort_2401_12963}
& New-task practice outcomes; foundation-model task proposals, affordances,
and safety checks.
& Generate practice or fleet experience, add it to a generalist mixture, and
fine-tune; AutoRT separates semantic task proposal from executable robot
collection.
& Safety rules and human/environment constraints gate collection.  The loop
does not autonomously attribute gains among proposer, data, objective, and
model changes. \\
\midrule
$\pi^{*}_{0.6}$/RECAP
~\cite{intelligence20250_2511_14759}
& Terminal success/failure labels, time-to-outcome values, autonomous trials,
and optional human interventions.
& Train a language-conditioned value function on accumulated experience and
extract an advantage-conditioned VLA; repeat collection and offline updates.
& Policy/value are restarted from the audited pretrained checkpoint at each
iteration in the reported implementation to limit drift.  Rewards and
interventions still require task-specific human infrastructure. \\
\midrule
DexFlyWheel
~\cite{zhu2025dexflywheel_2509_23829}
& Simulator reward and rollout success under new objects/configurations.
& One human seed bootstraps imitation learning; residual RL explores; successful
rollouts and geometric augmentation form the next iteration's dataset.
& Closed-loop simulation makes validity testable, but manual reward design and
the digital twin define the reachable distribution; it is not an open-world
failure-discovery system. \\
\midrule
SOP
~\cite{pan2026sop_2601_03044}
& Fresh on-policy failures/interventions and sliding-window online versus
offline training losses.
& A distributed actor fleet streams complete episodes to one learner; a
task-balanced sampler mixes static and online buffers and plug-in HG-DAgger or
RECAP updates are broadcast asynchronously.
& Uniform task weights and a clipped 0.2--0.8 online ratio preserve coverage;
checkpoints change only between episodes.  It still depends on human
interventions or task rewards and is tested over hours, not lifelong operation. \\
\midrule
LWD
~\cite{wang2026learning_2605_00416}
& Sparse terminal rewards over successes, failures, retries, partial progress,
and intervention segments from a 16-robot fleet.
& Distributional value learning uses heterogeneous offline/online replay;
critic gradients update a flow-based action expert through adjoint matching,
and the shared policy is repeatedly redeployed.
& A fixed offline reference, clipped double critic, mixed replay, and freezing
the VLM backbone online reduce instability.  Four-hour studies do not yet
establish long-term verifier or distribution stability. \\
\midrule
Robot-Powered Data Flywheels
~\cite{grannen2025robot_2511_19647}
& Deployment novelty and errors, with environmental/catalog signals that can
provide labels during use.
& A two-week library deployment collects task-relevant robot observations and
uses them to adapt an upstream visual foundation model, feeding improved
perception back to the robot system.
& Demonstrates that a flywheel may update perception rather than only the
action policy.  Site-specific labels and limited duration leave transfer and
long-run bias unresolved. \\
\midrule
DEED
~\cite{sis2026closing_2607_20345}
& Deployment experience and failures reveal gaps in an offline demonstration
set.
& Experience is converted into new or revised demonstrations, then used to
update the policy and the next data-collection decision.
& Makes demonstration evolution explicit, but causal promotion tests,
independent verifiers, and rollback remain necessary when several loop
components change together. \\
\bottomrule
\end{tabularx}
\end{table*}


The table also exposes what the term \emph{flywheel} can conceal. Simulator success, sparse task outcome, human takeover, and model uncertainty are not interchangeable feedback. A useful system description identifies the observed event, the unit derived from it, the objective that consumes the unit, and the control that prevents one round from damaging retained capabilities.

\subsection{Post-Training Changes Which States Matter}

Offline demonstrations are collected under an expert, operator, or scripted distribution. Data collected by the current policy expose the states and mistakes that the model itself produces. We call these records \emph{on-policy} when they arise from the policy being updated. Their value cannot be compared with offline data by hours alone. The Scalable Online Post-Training system (SOP) reports that three hours of targeted on-policy interaction closed more of the performance gap than an additional 80 hours of static demonstrations for one base model and task family~\cite{pan2026sop_2601_03044}. This is a within-paper result, not a universal exchange rate. It supports the narrower claim that targeted experience can have high marginal value near deployment.

The relevant preparation operation is not simply to add failures. A successful segment can be an imitation target. A valid prefix of a failed episode can support skill or value learning. An intervention supplies a takeover boundary and a corrective action. A near miss can train an outcome model, while an unsafe action may calibrate a safety verifier but should not silently enter behavior cloning. Ambiguous cases need uncertainty and review. RoboReward illustrates this role separation by revising outcome labels and clipping trajectories in time to construct negative and near-miss examples for a specialized reward vision--language model~\cite{lee2026roboreward_2601_00675}. Its real-robot study also shows that a general embodied reasoning model is not automatically a calibrated reward model for a target robot distribution.

Reinforcement learning and post-training increase the need for this preparation layer. A scalar reward does not identify which interval caused failure, whether an earlier prefix remains valid, or whether a low return reflects unsafe behavior, exploration, or label error. Outcome inference, temporal credit assignment, intervention parsing, and buffer routing are data decisions even when the later optimizer uses reinforcement learning. Reporting only the reward algorithm hides the evidence that determined what entered the update.

\subsection{Update Design Determines Stability}

RECAP, SOP, and Learning While Deploying (LWD) illustrate three update designs. RECAP combines demonstrations, autonomous rollouts, terminal outcomes, and optional human corrections. Its value model predicts time to success or a large failure penalty, and the policy favors actions estimated to improve over reference behavior~\cite{intelligence20250_2511_14759}. In its reported multi-round implementation, each value model and policy is fine-tuned from a common pretrained checkpoint instead of the previous round. This restart limits one path of accumulated drift.

SOP reduces the delay between failure and correction through distributed actors. It balances tasks, clips each task's online share to a bounded range, and changes robot checkpoints only between episodes~\cite{pan2026sop_2601_03044}. LWD replays offline and online experience, models a distribution of possible returns, and freezes the vision--language backbone during online updates while changing the action-producing module~\cite{wang2026learning_2605_00416}. These controls operate at different layers. Restarting controls model lineage, mixture bounds control the data distribution, and freezing modules limits the update path. Their effects should be tested separately rather than grouped as generic stabilization.

A multi-round evaluation needs at least three controls. Utility should be reported throughout the sequence of robot interaction instead of only at the final checkpoint. Previously learned and safety-critical capabilities need a fixed retention audit. The proposed data decision should also be compared with a volume- and compute-matched sample from the same collection round. Without this comparison, a stronger successor can reflect more interaction, more optimization, or a changed reward rather than the engine's routing decision.

Closed loops also create correlated error. The current policy controls which states are visited, a learned judge labels those states, and the resulting labels train the next policy. If the judge and generator share blind spots, repeated updates can make the missing region less visible. Stable systems therefore retain audited reference data, identify generated lineage, monitor the weakest task or embodiment group, and define deployment and rollback conditions before the next round begins.

\subsection{From Prescribed Engines to Preparation Agents}

Figure~\ref{fig:autonomy-ladder} organizes the progression by decision authority rather than publication date. At L0, people choose a fixed pipeline. L1 adapts an offline subset or mixture to a learner. L2 uses rollout failures or interventions for one or a few updates. An L3 engine repeatedly executes a predefined collect--verify--route--train controller. An L4 preparation agent additionally chooses which operator or diagnostic to run, estimates its cost and information value, and decides when to abstain, stop, or roll back. Existing systems demonstrate substantial components through L3. Describing them as fully autonomous L4 agents would overstate current evidence.

\begin{figure*}[t]
\centering
\begin{tikzpicture}[
  level/.style={draw, rounded corners=2pt, align=center, text width=2.88cm,
    minimum height=1.75cm, inner sep=4pt, font=\footnotesize},
  arr/.style={-{Latex[length=1.8mm]}, thick, draw=black!55},
  axis/.style={-{Latex[length=2mm]}, very thick, draw=surveyblue}
]
\node[level, fill=surveylightgray] (l0) {\textbf{L0: Static}\\fixed rules and one-shot mixture\\\emph{people choose the pipeline}};
\node[level, fill=surveylightblue, right=3mm of l0] (l1) {\textbf{L1: Learner-aware}\\loss, influence, or proxy selection\\\emph{offline adaptation}};
\node[level, fill=surveylightteal, right=3mm of l1] (l2) {\textbf{L2: Failure-driven}\\rollouts alter later training data\\\emph{one or a few rounds}};
\node[level, fill=surveylightorange, right=3mm of l2] (l3) {\textbf{L3: Data engine}\\repeated collect--verify--train loop\\\emph{fixed controller}};
\node[level, fill=surveylightrose, right=3mm of l3] (l4) {\textbf{L4: Preparation agent}\\chooses operators, evidence, and stopping\\\emph{adaptive controller}};
\draw[arr] (l0) -- (l1);
\draw[arr] (l1) -- (l2);
\draw[arr] (l2) -- (l3);
\draw[arr] (l3) -- (l4);
\draw[axis] ([yshift=-7mm]l0.south west) -- node[below, align=center, font=\small] {increasing decision authority and target relevance\\increasing feedback dependence, evidence cost, and safety risk} ([yshift=-7mm]l4.south east);
\end{tikzpicture}
\caption{Evolution of data preparation by decision authority. Existing systems demonstrate substantial components through L3. L4 additionally requires autonomous experiment choice, matched deployment tests, stopping, and rollback.}
\label{fig:autonomy-ladder}
\end{figure*}


An L4 agent is an experimental controller rather than an unrestricted data-generation model. It needs an operator library with explicit input and output schemas, versioned data and model lineage, calibrated evaluation, explicit budgets, and authority limits. At each round it can collect a missing case, revise an existing unit, change a training role, purchase stronger evidence, or stop. These choices return evidence at different delays and risks. A schema check is immediate, a controlled training run is delayed, and a real-robot test exposes physical cost.

The agent should optimize policy utility over the sequence of decisions rather than an immediate judge score. Every transformation needs a reversible record, and every deployment decision needs a matched comparison. A large language model (LLM) can plan operations or provide a semantic interface, but language fluency does not supply action validity or permission for risky exploration. The transition from L3 to L4 occurs only when the system can choose experiments, quantify uncertainty, preserve constraints, and replay the complete decision sequence.
